\section{告発の倫理：段階と証拠、改善への見込み}
\label{theory:ethics:whistle_conditions}
深刻な危害が見込まれるとき、内部で懸念を伝え有効な是正を求める。
上司が応答しないという一事だけで、他の内部経路も調べず直ちに外部公表へ進むとはしない。
客観的な証拠、危害を伝える相手、公表による改善の見込みと個人に及ぶ危険を考える。
これは倫理的な意思決定を整理する枠組みで、法定の公益通報保護要件そのものではない。\cite{muroran_ethics_cases}

令和元年度再試験の比較では、証拠がなくても外部へ知らせる義務が必ずあるという肢を否定する。
一方、危害が急迫している実際の事態に、全ての内部段階を必ず順番に待つという手順を命じるものでもない。
緊急性と有効な通報先、必要な開示範囲、現在の法定経路を調べる。
原著全章を確認した解説ではないため、倫理モデルの解釈は査読事項として保持する。
